\renewcommand{\theposttitle}{I.D.E.A. Newsletter 9: Storytelling and Comfort Zones}
\renewcommand{\thepostdate}{2022-07-22}
\vspace*{\fill}
\begin{center}
\includegraphics[width=\textwidth, height=0.24\textheight, keepaspectratio]{C:/Users/James/Documents/GitHub/VisualIdeasBook/data/images/998ae232699231de89393e811985d724.png}
\end{center}
\vspace{5mm}
\begin{center}
\includegraphics[width=\textwidth, height=0.24\textheight, keepaspectratio]{C:/Users/James/Documents/GitHub/VisualIdeasBook/data/images/1707e8138ed61b9220cdd548236b3de6.png}
\end{center}
\vspace{2mm}
\begin{center}
\begin{minipage}{0.92\textwidth}
\small
The very last lines of "The Practice" by Seth Godin reads:

"The magic is that there is no magic. Start where you're at. Don't Stop."

The top of the Maslow's Hierarchy of Needs pyramid should be changed from "Self-Actualization" to "Continuous Progress."

There's a reason why we begin most movies with our hero in an Ordinary World.

This is the comfort zone.

Soon after our hero's introduced, there's a Call to Action, followed swiftly by the hero's Refusal of the Call.

Why?

Because few people intentionally step into the unknown and into discomfort.

They often have to be pushed into it.

The tragic truth of growth is that it's something we recognize only in hindsight and not as it's happening.

But growth requires us to step into the fray and trust we'll get through to the other side stronger for it.

A flower only grows when it reaches for the sun.
\end{minipage}
\end{center}
\vspace*{\fill}